\documentclass[12pt,a4paper]{article}

% -------------------- Packages --------------------
\usepackage[utf8]{inputenc}
\usepackage{geometry}
\geometry{left=1in, right=1in, top=1in, bottom=1in}

\usepackage{graphicx}
\usepackage{setspace}
\usepackage{titlesec}
\usepackage{fancyhdr}
\usepackage{listings}
\usepackage{xcolor}
\usepackage{float}
\usepackage{hyperref}
\usepackage[utf8]{inputenc} 
\usepackage[margin=1in]{geometry} 
\usepackage{framed}
\usepackage{listings} 
\usepackage{xcolor} 
\usepackage{hyperref} 
\usepackage{fancyhdr} 
\usepackage{titlesec} 
\usepackage{float} 

\let\oldverbatim\verbatim
\let\endoldverbatim\endverbatim

\renewenvironment{verbatim}
{
    \vspace{0.5em}
    \begin{framed}
    \noindent\oldverbatim
}
{
    \endoldverbatim
    \end{framed}
    \vspace{0.5em}
}

% Code listing style 
\definecolor{codegreen}{rgb}{0,0.6,0} 
\definecolor{codegray}{rgb}{0.5,0.5,0.5} 
\definecolor{codepurple}{rgb}{0.58,0,0.82} 
\definecolor{backcolour}{rgb}{0.95,0.95,0.92}

\lstdefinestyle{mystyle}{ backgroundcolor=\color{backcolour}, commentstyle=\color{codegreen}, keywordstyle=\color{blue}, numberstyle=\tiny\color{codegray}, stringstyle=\color{codepurple}, basicstyle=\ttfamily\footnotesize, breakatwhitespace=false, breaklines=true, captionpos=b, keepspaces=true, numbers=left, numbersep=5pt, showspaces=false, showstringspaces=false, showtabs=false, tabsize=2 } \lstset{style=mystyle}

% -------------------- Line Spacing --------------------
\onehalfspacing

% -------------------- Header / Footer --------------------
\pagestyle{fancy}
\fancyhf{}
\rhead{Web Engineering Lab}
\lhead{Lab 14: Concurrency in Java}
\cfoot{\thepage}

% -------------------- Section Formatting --------------------
\titleformat{\section}
  {\normalfont\bfseries\uppercase}{\thesection.}{1em}{}

\titleformat{\subsection}
  {\normalfont\bfseries}{\thesubsection}{1em}{}

% -------------------- Code Style (Simple Academic) --------------------
\lstset{
    basicstyle=\ttfamily\footnotesize,
    numbers=left,
    numbersep=5pt,
    frame=single,
    breaklines=true,
    tabsize=2,
    showstringspaces=false
}

% -------------------- Document Start --------------------
\begin{document}


\begin{titlepage}
\centering

\vspace*{0.5in}

\includegraphics[width=0.22\textwidth]{image.png}\par
\vspace{0.4in}

\textbf{\Large National University of Sciences \& Technology}\\
\textbf{\large School of Electrical Engineering and Computer Science (SEECS)}\\[0.3in]

\textbf{\large Web Engineering Laboratory}\\[0.5in]

\textbf{\Large Lab 14: Concurrency in Java}\\[0.4in]

\textbf{Submitted By}\\[0.15in]
\textbf{Mustafa Hamad}\\
\textbf{CMS: 455095}\\
\textbf{Section: SE-14A}\\[0.4in]

\textbf{Submitted To}\\[0.15in]
\textbf{Sir Aftab Farooq}\\[0.4in]

\textbf{Department of Computer Science}\\
\textbf{NUST, Islamabad}\\[0.4in]

\textbf{December 15, 2025}

\end{titlepage}

\newpage

\tableofcontents
\newpage

\section{Introduction}
This lab focuses on understanding and implementing concurrency concepts in Java. Concurrency allows multiple threads to execute simultaneously, improving application performance and responsiveness. However, it introduces challenges such as race conditions, deadlocks, and data inconsistency, which require careful handling through synchronization mechanisms.

\subsection{Objectives}
By the end of this lab, students will be able to:
\begin{itemize}
    \item Create and execute multiple threads in Java
    \item Implement thread synchronization to avoid race conditions
    \item Use thread-safe data structures for concurrent operations
    \item Simulate real-world concurrent scenarios (bank transactions)
\end{itemize}

\section{Task 1: Introduction to Multithreading}

\subsection{Objective}
To create and execute multiple threads in Java and observe concurrent execution behavior.

\subsection{Implementation}
Two threads are created using the \texttt{Runnable} interface:
\begin{itemize}
    \item \textbf{NumberThread}: Prints numbers from 1 to 10
    \item \textbf{SquareThread}: Prints squares of numbers from 1 to 10
\end{itemize}

\newpage
\subsection{Key Code}
\begin{lstlisting}[language=Java, caption=Thread Creation and Execution]
// Create threads using Runnable interface
Thread numberThread = new Thread(new NumberPrinter(), "NumberThread");
Thread squareThread = new Thread(new SquarePrinter(), "SquareThread");

// Start both threads
numberThread.start();
squareThread.start();

// Wait for completion
numberThread.join();
squareThread.join();
\end{lstlisting}

\begin{lstlisting}[language=Java, caption=NumberPrinter Runnable]
class NumberPrinter implements Runnable {
    @Override
    public void run() {
        for (int i = 1; i <= 10; i++) {
            System.out.println("[NumberThread] Number: " + i);
            Thread.sleep(100);
        }
    }
}
\end{lstlisting}

\subsection{Output}
\begin{verbatim}
Starting both threads concurrently...

[NumberThread] Number: 1
[SquareThread] Square of 1 = 1
[NumberThread] Number: 2
[SquareThread] Square of 2 = 4
[NumberThread] Number: 3
[SquareThread] Square of 3 = 9
...
Both threads have completed execution.
\end{verbatim}

\subsection{Observations}
The output demonstrates interleaved execution where both threads run concurrently. The order of output may vary between runs due to the non-deterministic nature of thread scheduling.

\newpage
\section{Task 2: Thread Synchronization}

\subsection{Objective}
To implement thread synchronization to avoid race conditions when multiple threads access shared resources.

\subsection{Problem Statement}
Three threads increment a shared counter 100 times each. Without synchronization, the final value would be unpredictable due to race conditions. With proper synchronization, the final value is guaranteed to be 300.

\subsection{Key Code}
\begin{lstlisting}[language=Java, caption=Synchronized Counter Class]
class SharedCounter {
    private int count = 0;

    // Synchronized method ensures only one thread can execute at a time
    public synchronized void increment() {
        count++;
    }

    public int getCount() {
        return count;
    }
}
\end{lstlisting}

\begin{lstlisting}[language=Java, caption=Counter Incrementer Thread]
class CounterIncrementer implements Runnable {
    private final SharedCounter counter;

    public CounterIncrementer(SharedCounter counter) {
        this.counter = counter;
    }

    @Override
    public void run() {
        for (int i = 0; i < 100; i++) {
            counter.increment();
        }
    }
}
\end{lstlisting}

\subsection{Output}
\begin{verbatim}
Starting three threads to increment a shared counter...

[Thread-1] Completed 100 increments.
[Thread-2] Completed 100 increments.
[Thread-3] Completed 100 increments.

All threads have completed.
Final Counter Value: 300
Expected Value: 300
SUCCESS: Synchronization worked correctly!
\end{verbatim}

\subsection{Key Concepts}
\begin{itemize}
    \item \textbf{Race Condition}: Occurs when multiple threads access shared data without synchronization
    \item \textbf{synchronized keyword}: Ensures mutual exclusion - only one thread can execute the synchronized method at a time
    \item \textbf{Thread Safety}: The property of code that behaves correctly when accessed by multiple threads
\end{itemize}

\newpage
\section{Task 3: Concurrent Data Structures}

\subsection{Objective}
To implement and use thread-safe data structures for concurrent operations without explicit synchronization.

\subsection{Data Structures Used}
\begin{itemize}
    \item \textbf{CopyOnWriteArrayList}: A thread-safe variant of ArrayList where all mutative operations are implemented by making a fresh copy of the underlying array
    \item \textbf{ConcurrentHashMap}: A thread-safe hash map that allows concurrent read and write operations
\end{itemize}

\subsection{Key Code}
\begin{lstlisting}[language=Java, caption=Declaring Thread-Safe Collections]
import java.util.concurrent.ConcurrentHashMap;
import java.util.concurrent.CopyOnWriteArrayList;

// Thread-safe list
private static final List<Integer> sharedList = new CopyOnWriteArrayList<>();

// Thread-safe map
private static final Map<String, Integer> sharedMap = new ConcurrentHashMap<>();
\end{lstlisting}

\begin{lstlisting}[language=Java, caption=Concurrent List Writer]
class ListWriter implements Runnable {
    private final List<Integer> list;
    
    @Override
    public void run() {
        for (int i = start; i <= end; i++) {
            list.add(i); // Thread-safe operation
        }
    }
}
\end{lstlisting}

\subsection{Output}
\begin{verbatim}
=== Demonstrating CopyOnWriteArrayList ===

[ListWriter-1] Added: 1
[ListReader-1] Reading list: [1]
[ListWriter-2] Added: 6
[ListWriter-1] Added: 2
...
Final List Contents: [1, 6, 2, 7, 3, 8, 4, 9, 5, 10]
List Size: 10

=== Demonstrating ConcurrentHashMap ===

[MapWriter-1] Put: A1 -> 10
[MapReader-1] Reading map: {A1=10}
...
Final Map Contents: {A1=10, A2=20, A3=30, B4=40, B5=50, B6=60}
Map Size: 6
\end{verbatim}

\subsection{Advantages}
\begin{itemize}
    \item No explicit synchronization needed
    \item Better performance under high contention
    \item Built-in thread safety guarantees
\end{itemize}

\newpage
\section{Task 4: Bank Transaction System Simulation}

\subsection{Objective}
To simulate a bank account system with multiple clients performing concurrent deposits and withdrawals while maintaining data integrity.

\subsection{Implementation Details}
\begin{itemize}
    \item 5 client threads perform random transactions
    \item \textbf{AtomicInteger} is used for thread-safe balance operations
    \item \textbf{synchronized} methods prevent race conditions in withdrawals
\end{itemize}

\subsection{Key Code}
\begin{lstlisting}[language=Java, caption=Thread-Safe Bank Account]
class BankAccount {
    private final AtomicInteger balance;
    private final AtomicInteger totalDeposits = new AtomicInteger(0);
    private final AtomicInteger totalWithdrawals = new AtomicInteger(0);

    public void deposit(String clientName, int amount) {
        int newBalance = balance.addAndGet(amount);
        totalDeposits.addAndGet(amount);
    }

    public synchronized boolean withdraw(String clientName, int amount) {
        if (balance.get() >= amount) {
            int newBalance = balance.addAndGet(-amount);
            totalWithdrawals.addAndGet(amount);
            return true;
        }
        return false; // Insufficient funds
    }
}
\end{lstlisting}
\newpage
\begin{lstlisting}[language=Java, caption=Bank Client Thread]
class BankClient implements Runnable {
    @Override
    public void run() {
        for (int i = 0; i < 5; i++) {
            int amount = random.nextInt(100) + 10;
            if (random.nextBoolean()) {
                account.deposit(clientName, amount);
            } else {
                account.withdraw(clientName, amount);
            }
        }
    }
}
\end{lstlisting}

\subsection{Output}
\begin{verbatim}
=== Bank Transaction System Simulation ===
Initial Balance: $1000
Starting 5 client threads...

[Alice] DEPOSIT: $75 | Balance: $1075
[Bob] WITHDRAW: $42 | Balance: $1033
[Charlie] DEPOSIT: $89 | Balance: $1122
[Diana] WITHDRAW: $56 | Balance: $1066
[Eve] DEPOSIT: $31 | Balance: $1097
...

=== All Transactions Completed ===
Final Balance: $1156
Total Deposits: $432
Total Withdrawals: $276
Expected Final Balance: $1156
\end{verbatim}

\subsection{Key Concepts}
\begin{itemize}
    \item \textbf{AtomicInteger}: Provides atomic (indivisible) operations on integers without explicit synchronization
    \item \textbf{addAndGet()}: Atomically adds the given value and returns the updated value
    \item \textbf{Synchronized withdrawal}: Prevents overdraft by checking balance before withdrawal atomically
\end{itemize}

\newpage
\section{Conclusion}
This lab provided hands-on experience with Java concurrency concepts:

\begin{enumerate}
    \item \textbf{Task 1}: Demonstrated basic multithreading using \texttt{Thread} and \texttt{Runnable}
    \item \textbf{Task 2}: Showed how \texttt{synchronized} keyword prevents race conditions
    \item \textbf{Task 3}: Introduced thread-safe collections (\texttt{CopyOnWriteArrayList}, \texttt{ConcurrentHashMap})
    \item \textbf{Task 4}: Applied concurrency concepts to a real-world banking scenario using \texttt{AtomicInteger}
\end{enumerate}

\subsection{Key Takeaways}
\begin{itemize}
    \item Concurrent programming requires careful consideration of shared resources
    \item The \texttt{synchronized} keyword provides basic thread safety but may impact performance
    \item Java's concurrent utilities (\texttt{java.util.concurrent}) offer optimized thread-safe alternatives
    \item Atomic classes provide lock-free thread safety for simple operations
\end{itemize}

\section{References}
\begin{itemize}
    \item Oracle Java Documentation: Concurrency Tutorial
    \item Java SE API: java.util.concurrent package
    \item "Java Concurrency in Practice" by Brian Goetz
\end{itemize}

\end{document}
